\documentclass[11pt]{article}

\usepackage[margin=1in]{geometry}
\usepackage{booktabs}
\usepackage{graphicx}
\usepackage{hyperref}
\usepackage{array}
\usepackage{longtable}
\usepackage{xcolor}
\usepackage{amsmath}

\hypersetup{
  colorlinks=true,
  linkcolor=blue,
  urlcolor=blue
}

\title{CSCI-360 Machine Learning Final Competition 26s \\ CSI500 Stock Selection Report}
\author{Mingqian Yang}
\date{11 May 2026}

\begin{document}
\maketitle

\noindent\textbf{Path note.} All file paths and commands referenced in this
report are relative to the cleaned project repository available at
\url{https://github.com/TimmoTaffy/ml-competition-sp26}. This note matters
because an earlier Phase2 code ZIP was submitted before the documentation
cleanup and reproducibility edits, and the current Gradescope report channel
does not accept a code ZIP upload.

\section*{Executive Summary}

The final system trains XGBoost regressors from scratch to rank CSI500 stocks
by expected five-trading-day excess return over the CSI500 index. The Phase1
submission used the strongest single candidate, \texttt{stable\_compact\_current}.
After the Phase1 window became public, Phase2 used a more robust
portfolio-level blend of \texttt{stable\_compact\_current} and
\texttt{stable\_compact\_capacity}: \(0.45\) current plus \(0.55\) capacity,
truncated to 44 holdings.

The workflow does not reserve a separate static test set for final model selection. This is
intentional: the task is a live portfolio competition with very few recent
trading days, and the final judge is the next unseen holding window. Instead,
the workflow uses strict time-causal live-style walk-forward validation for
model and portfolio selection, then reports the two course live windows as
external out-of-sample checks. The Phase1 submitted portfolio earned
\(+4.834\%\) excess return from May 6--8, compared with \(+2.693\%\) for the
instructor XGBoost baseline regenerated under the same as-of convention.

\clearpage

\section{Factors}

The feature design started from the view that most single short-horizon market
features are noisy. Instead of relying on one hand-coded rule, the project
constructed a large but economically grouped feature pool and selected features
by rolling validation.

The report lists factor groups and representative examples rather than every
raw column. The complete reproducible feature universe is in
\texttt{data/final\_feature\_columns.txt}; the exact retained lists are in
\texttt{experiments/research\_families/fine\_tune\_top12/feature\_sets/}; and
the feature engineering notes are in
\texttt{docs/FEATURES\_AND\_SELECTION.md}.

\begin{longtable}{p{0.24\linewidth}p{0.33\linewidth}p{0.34\linewidth}}
\toprule
Group & Examples & Rationale \\
\midrule
\endfirsthead
\toprule
Group & Examples & Rationale \\
\midrule
\endhead
Technical and momentum & 1/2/3/5/10/20/60/120-day returns, volatility, RSI, moving-average distance & Captures short-horizon continuation and reversal patterns. \\
CSI500-relative & excess returns versus CSI500, beta, residual momentum, relative volatility & Aligns the prediction target with the live excess-return objective. \\
Bollinger bands & band z-score, position, width, lower/upper proximity, near/breakout flags & Tests mean-reversion and breakout effects without hard-coding direction. \\
Liquidity and risk & halt flags, zero-volume days, missing days, turnover, ATR, drawdown & Avoids wasting weight on untradeable or high-tail-risk names. \\
Industry-relative & Shenwan industry returns and stock-minus-industry ranks & Separates stock selection from sector rotation. \\
Valuation and quality & PE/PB, book-to-price, ROE, ROA, margins, 3-year revenue/profit growth, cash-flow quality & Adds slower-moving economic information beyond price noise. \\
Market regime & broad indices, ETFs, sector ETF proxies, crude/gold/global proxies, CSI500 quality-growth index & Allows the model to react to style and macro regime shifts. \\
\bottomrule
\end{longtable}

AH-premium features were removed because CSI500 dual-listed coverage was too
low for a stable cross-sectional signal. Feature selection was profile-specific:
the general model and short-history model used separate selected lists. The
candidate ranking families \texttt{balanced}, \texttt{stable}, \texttt{recent},
\texttt{economic}, and \texttt{short\_aggressive} changed the feature-selection
objective, then all candidates were compared by the same portfolio validation
standard.

\subsection{Feature Selection Process}

The selection pipeline was designed to avoid trusting one model's feature
importance. For each profile, features were first tested as single factors over
rolling historical windows. The single-factor checks included rank IC, top
decile forward excess return, top-minus-bottom spread, hit rate, worst-window
behavior, missingness, and cross-window stability.

Next, grouped model tests checked whether economically coherent feature groups
improved a portfolio-like objective. This matters because a feature with low
standalone IC can still help in combination, while a noisy feature can look good
in one window by chance. The combined ranking score rewarded stable positive
rank IC, positive top-decile excess, good downside behavior, and useful group
ablation results, while penalizing unstable or sparse signals.

\begin{table}[h]
\centering
\begin{tabular}{p{0.28\linewidth}p{0.60\linewidth}}
\toprule
Ranking family & Reasoning \\
\midrule
\texttt{balanced} & Broad default score across IC, top-decile return, and stability. \\
\texttt{stable} & Emphasizes lower missingness, downside control, and stable signs. \\
\texttt{recent} & Gives more weight to recent windows, useful for regime shifts. \\
\texttt{economic} & Gives more weight to quality, valuation, and fundamental groups. \\
\texttt{short\_aggressive} & Emphasizes short-horizon technical and recent signals. \\
\bottomrule
\end{tabular}
\end{table}

The reason for keeping multiple ranking families is that the feature-ranking
step interacts with model capacity and portfolio construction. Instead of
declaring one feature ranking ``correct'', I let several economically defensible
rankings compete under the same live-style portfolio validation.

Concrete feature-selection inputs are in
\path{configs/feature_selection_general.json} and
\path{configs/feature_selection_short_history.json}. The ranking-family logic is
implemented in \texttt{tune\_research\_families.py}. The selected feature lists
for retained candidates are under \path{experiments/.../feature_sets/}; the
complete path map is in \path{docs/MODEL_SELECTION_LINEAGE.md}.

\clearpage

\section{Models}

Each source model is an \texttt{XGBRegressor} trained from scratch. A training
row is one stock on one date, and the model predicts \texttt{target\_excess\_5d}.
The final structure has two profiles:

\begin{table}[h]
\centering
\begin{tabular}{p{0.25\linewidth}p{0.64\linewidth}}
\toprule
Profile & Role \\
\midrule
\texttt{general} & Broader cross-sectional signal with more history and more features. \\
\texttt{short\_history} & Recent technical/regime signal with shorter training lookback. \\
\bottomrule
\end{tabular}
\end{table}

Predictions are converted to cross-sectional percentile ranks before blending:
\[
\text{combined score}
= w_g \cdot \text{rank}(\hat y_g)
+ w_s \cdot \text{rank}(\hat y_s).
\]
Using ranks avoids mixing incompatible raw XGBoost score scales. The portfolio
builder sorts by combined score, applies the risk filter, selects \texttt{top\_k},
assigns equal or rank weights, caps concentration, and renormalizes to one.

The general/short split was kept because the two profiles answer different
questions. The general model uses broader history and is better suited for
stable cross-sectional relationships. The short-history model uses a shorter
lookback and is better suited for recent technical or regime effects. Their raw
predictions are not directly comparable, so the ensemble uses percentile ranks
instead of raw scores. The weight between them is not fixed by hand; it is a
portfolio hyperparameter selected during validation.

The default profile configs are \path{configs/model_general.json} and
\path{configs/model_short_history.json}. The final current variants are
\texttt{general\_top80\_shallow} and
\texttt{short\_history\_top60\_recent\_train}; the final capacity variants are
\texttt{general\_top80\_base} and
\texttt{short\_history\_top20\_deep\_interaction}. Exact JSON paths and selected
feature-list paths are listed in \path{docs/HYPERPARAMETERS.md} and
\path{docs/PHASE_REPRODUCTION.md}. Final ensemble/portfolio settings are in
\path{configs/ensemble_phase1.json}, \path{configs/ensemble_stable_compact_capacity.json},
and the Phase2 blend config in \path{experiments/portfolio_blend_current_capacity/}.

\subsection{Hyperparameter And Candidate Search}

Hyperparameters were selected through a staged workflow:
\begin{enumerate}
  \item Build the feature matrix from public data with announcement-date-aware fundamentals.
  \item Run rolling feature selection for general and short-history profiles.
  \item Search structured research families: ranking family, feature-count pair, and XGB preset.
  \item Retain a frozen top12 candidate set.
  \item Re-evaluate retained candidates with strict live-style walk-forward validation.
  \item Fine-tune only portfolio construction on 3/7/13 recent non-overlapping 5-day windows.
  \item For Phase2, tune a portfolio-level current/capacity blend over alpha and \texttt{top\_k}.
\end{enumerate}

The structured research search crossed:
\[
5 \text{ feature-ranking families}
\times 4 \text{ feature-count pairs}
\times 3 \text{ XGBoost presets}
= 60 \text{ candidate families}.
\]
The feature-count pairs were \texttt{compact}, \texttt{medium}, \texttt{wide},
and \texttt{extra\_wide}. The XGBoost presets were \texttt{conservative},
\texttt{current}, and \texttt{capacity}. \texttt{conservative} used
shallower/more regularized trees; \texttt{current} was close to a reasonable
baseline tree shape; and \texttt{capacity} allowed more interactions. The
search first narrowed 60 candidates to a smaller pool, then froze 12 candidates
for strict live-style validation and portfolio tuning.

The 60-candidate construction is documented in
\path{docs/MODEL_SELECTION_LINEAGE.md} and implemented in
\texttt{tune\_research\_families.py}. The retained top33/top12 audit trail is in
\path{candidate_family_lookup.csv} under the fine-tune experiment directory; the
frozen top12 candidate directories are under \path{experiments/.../candidates/}.
The readable top12 rerun guide is \path{docs/TOP12_REPRODUCTION.md}.

The final candidate comparison did not optimize standalone IC. It optimized a
portfolio-oriented validation score. For each historical window, the system
retrained from scratch, created a legal portfolio, and measured realized excess
return. Candidate scoring combined mean excess, worst-window excess, and hit
rate, because the live competition rewards realized portfolio excess rather
than prediction correlation alone.

Live-style validation was also used to compare the final Phase2 source
candidates. For each historical 5-day window, the system retrained from scratch
at the entry date using only labels known then, built a legal portfolio, and
scored the following window. On the latest 13-window run ending 2026-05-08:

\begin{table}[h]
\centering
\begin{tabular}{lrrrr}
\toprule
Candidate & Mean excess & Hit rate & Worst window & Live score \\
\midrule
\texttt{stable\_compact\_capacity} & \(+0.881\%\) & \(76.9\%\) & \(-2.265\%\) & 0.795 \\
\texttt{stable\_compact\_current} & \(+1.211\%\) & \(61.5\%\) & \(-3.734\%\) & 0.711 \\
\bottomrule
\end{tabular}
\end{table}

The current candidate had higher mean excess but worse drawdown stability. The
capacity candidate had lower mean but better hit rate and worst-window control.
This tradeoff motivated Phase2's portfolio-level blend. Complete top12
rankings, including candidates not used in the final Phase2 blend, are recorded
in \path{docs/TOP12_REPRODUCTION.md} and
\path{experiments/.../live_portfolio_fine_tune_3_7_13/leaderboard.csv}.

\begin{table}[h]
\centering
\begin{tabular}{p{0.28\linewidth}p{0.28\linewidth}p{0.34\linewidth}}
\toprule
Item & Phase1 & Phase2 \\
\midrule
Source candidate(s) & \texttt{stable\_compact\_current} & \texttt{stable\_compact\_current}, \texttt{stable\_compact\_capacity} \\
Current variants & \texttt{general\_top80\_shallow}, \texttt{short\_history\_top60\_recent\_train} & same \\
Capacity variants & not used & \texttt{general\_top80\_base}, \texttt{short\_history\_top20\_deep\_interaction} \\
Current score weight & 0.55 / 0.45 & inside source portfolio \\
Capacity score weight & not used & 0.978 / 0.022 \\
Final portfolio & 31-name rank-weighted & 0.45 current + 0.55 capacity, top 44 \\
Risk filter & enabled & enabled in both source portfolios \\
\bottomrule
\end{tabular}
\end{table}

\subsection{Phase1 Selection}

Phase1 was selected before the first live window. The strongest retained
candidate at that point was \texttt{stable\_compact\_current}. The design
process was:

\begin{enumerate}
  \item Generate 60 structured candidates from five feature-ranking families, four feature-count pairs, and three XGBoost presets.
  \item Use the early broad/fine research stage to reduce the search space to a top33 pool, then freeze 12 candidate families for final audit.
  \item Re-evaluate those 12 under live-style walk-forward validation.
  \item Fine-tune only portfolio construction on the retained candidates: general/short rank weight, \texttt{top\_k}, internal cap, risk filter, and equal/rank weighting.
  \item Run a robustness audit around the best rounded parameters.
\end{enumerate}

The final Phase1 choice was \texttt{stable\_compact\_current}. It combined
\texttt{general\_top80\_shallow} and
\texttt{short\_history\_top60\_recent\_train}. The fine-tuned winner used
approximately \(0.55/0.45\) general/short rank blending and a 31-stock
rank-weighted portfolio. The exact validation winner had
\(\text{general\_weight}=0.5465\) and \(\text{internal\_cap}=0.031758\); these
were rounded to \(0.55\) and \(0.032\) because the robustness audit showed
negligible validation score loss.

\subsection{Phase2 Selection}

Phase2 was selected after the May 6--8 Phase1 window became public. At that
point, \texttt{stable\_compact\_current} still had high mean excess but also
showed larger worst-window risk and parameter sensitivity.
\texttt{stable\_compact\_capacity} had lower mean excess but better robustness
in the latest live-style validation.

The key design decision for Phase2 was to blend portfolios, not raw model
scores. This preserves each source candidate's separately tuned legal
portfolio construction and avoids assuming that raw XGBoost scores from
different candidates are on the same scale. The Phase2 tuner compared pure
current, pure capacity, and portfolio blends on the same 3/7/13 live-style
validation anchors.

The final grid searched current weight \(\alpha\) and final \texttt{top\_k}.
Both the active robust score and the older 3/7/13 score selected
\texttt{final\_portfolio\_a0.45\_k44}, i.e. \(0.45\) current plus \(0.55\)
capacity with 44 final holdings. This choice kept part of the high-mean current
portfolio while reducing dependence on its weaker worst-window behavior.

\clearpage

\section{Results}

\subsection{Phase2 Blend Evidence}

The Phase2 alpha/top-K tuner selected \texttt{final\_portfolio\_a0.45\_k44}.
Both the active robust score and the older 3/7/13 score selected the same blend.
The table below compares the final blend to the two pure source candidates on
the same latest 13-window live-style validation run.

\begin{table}[h]
\centering
\begin{tabular}{lrrrr}
\toprule
Strategy & Robust score & 13w mean excess & 13w worst & 13w hit rate \\
\midrule
Phase2 blend, alpha 0.45, top 44 & 0.558 & \(+0.908\%\) & \(-2.926\%\) & \(53.8\%\) \\
Pure capacity & 0.481 & \(+0.768\%\) & \(-2.458\%\) & \(55.8\%\) \\
Pure current & 0.300 & \(+0.985\%\) & \(-3.734\%\) & \(57.7\%\) \\
\bottomrule
\end{tabular}
\end{table}

The blend sacrifices some pure-current mean return to reduce reliance on the
more parameter-sensitive current portfolio. It also keeps more upside than pure
capacity. This is the main empirical reason for selecting the Phase2 blend
instead of either pure candidate.

\subsection{External Check Against Baseline}

The Phase1 portfolio was generated before the May 6--8 prices were known. The
same realized window can therefore be used as an external post-submission check
of the Phase1 route. It is not used as a static final test for Phase2, because
its result was known before Phase2 tuning.

\begin{table}[h]
\centering
\begin{tabular}{lrrrr}
\toprule
Model & Holdings & Portfolio return & CSI500 return & Excess return \\
\midrule
Instructor baseline & 50 & \(+6.821\%\) & \(+4.128\%\) & \(+2.693\%\) \\
Phase1 model & 31 & \(+8.962\%\) & \(+4.128\%\) & \(+4.834\%\) \\
\bottomrule
\end{tabular}
\end{table}

The Phase1 model exceeded the instructor baseline by \(+2.141\) percentage
points of excess return on this external realized check. The exact
train/validation/test interpretation of this comparison is documented in the
Self-test section.

\clearpage

\section{Analysis}

\textbf{What worked.} Benchmark-relative labels and features worked better than
raw return prediction because the live score is excess return over CSI500.
Live-style walk-forward validation was more realistic than a single fixed
validation block because each window retrains as if it were a real submission
date. Rank-based score blending made the general and short-history XGBoost
outputs comparable, while liquidity/risk filters and capped portfolios reduced
implementation risk. The Phase2 current/capacity portfolio blend worked because
it kept part of the high-mean current candidate while reducing dependence on
its weaker worst-window behavior.

\textbf{What did not work.} Standalone IC was too noisy to select portfolios by
itself. Short-history-only configurations sometimes looked strong in isolated
windows but were unstable. AH-premium data had too little CSI500 coverage and
was removed. Pure current had attractive mean excess but larger worst-window
losses and higher parameter sensitivity, so Phase2 did not use it alone.

\textbf{Why the evaluation is still credible.} The main risk is overfitting to
recent historical windows. This is why the report avoids presenting a repeatedly
inspected fixed block as a clean test set. Instead, validation evidence is
separated from external realized evidence: validation chooses the model, Phase1
is reported as a realized post-submission check against the instructor baseline,
and Phase2 was the true unseen course evaluation for the final submitted
portfolio.

\clearpage

\section{Self-test}

\subsection{Train/Validation/Test Protocol}

The supervised target is \texttt{target\_excess\_5d}: a stock's future
five-trading-day return minus the future five-trading-day CSI500 return.
Because the label looks five trading days forward, every training and
validation step uses a five-day cutoff or embargo.

\begin{table}[h]
\centering
\begin{tabular}{p{0.23\linewidth}p{0.68\linewidth}}
\toprule
Layer & Definition \\
\midrule
Training & For each as-of date, train only on rows whose 5-day forward labels are fully known by that date. For example, a 2026-04-30 prediction row can only use labels through 2026-04-23. \\
Internal validation & Time-ordered validation blocks before the prediction date, with a five-trading-day embargo, are used for XGBoost early stopping and diagnostics. \\
Model/portfolio validation & Candidate comparison uses live-style walk-forward validation: retrain from scratch at each historical entry date, build a legal portfolio, and score the following 5-day window. \\
Held-out test / self-test & The Phase1 live window, 2026-05-06 to 2026-05-08, is used as the reported test set against the instructor baseline. It was not known when the Phase1 portfolio was generated. \\
No reused static test block & After a realized window has been inspected, it is no longer treated as a clean test for later Phase2 tuning. Phase2's true test was the unseen course live window after submission. \\
\bottomrule
\end{tabular}
\end{table}

This is a train/validation/test split adapted to a sequential portfolio
problem. Validation is not a random holdout; it is a set of historical
live-style simulations. The test set reported here is the first submitted
course window after it became realized. The protocol avoids the main leakage
risk in this task: a 5-day forward label whose future return window overlaps
the validation or live holding window. After hyperparameters are frozen, the
final live model is refit on all labels known at the submission date.

\subsection{Concrete Self-test Split}

The Phase1 portfolio was generated before the May 6--8 prices were known. This
is therefore the report's concrete self-test against the instructor baseline:
both portfolios are evaluated on the same realized out-of-sample holding
window, and the model pipeline was fixed before that window's labels were
known.

\begin{table}[h]
\centering
\begin{tabular}{p{0.16\linewidth}p{0.25\linewidth}p{0.23\linewidth}p{0.25\linewidth}}
\toprule
Split & Dates / cutoff & Use & Leakage control \\
\midrule
Training & As-of 2026-04-30; labels known through 2026-04-23 & Train Phase1 model from scratch & No May 6--8 price or label information is available \\
Validation & Historical pre-2026-04-30 live-style windows with 5-day embargo & Feature/model/portfolio selection & Each simulated entry date trains only on labels known then \\
Test & 2026-05-06 to 2026-05-08 & Report held-out performance versus baseline & Scored only after Phase1 submission \\
Baseline & Rebuilt as-of 2026-04-30 & Same test-window comparison & Same price window and benchmark return \\
\bottomrule
\end{tabular}
\end{table}

\begin{table}[h]
\centering
\begin{tabular}{lrrrr}
\toprule
Model & Holdings & Portfolio return & CSI500 return & Excess return \\
\midrule
Instructor baseline & 50 & \(+6.821\%\) & \(+4.128\%\) & \(+2.693\%\) \\
Phase1 model & 31 & \(+8.962\%\) & \(+4.128\%\) & \(+4.834\%\) \\
\bottomrule
\end{tabular}
\end{table}

The Phase1 model exceeded the instructor baseline by \(+2.141\) percentage
points of excess return on this held-out external check.

\subsection{Data And Leakage Controls}

The data pipeline uses public market data, public index/ETF proxies, Shenwan
industry data, valuation data, and announcement-date-aware fundamentals. The
important leakage controls are:

\begin{itemize}
  \item Training labels are cut off by the 5-trading-day horizon, so a prediction row never trains on labels whose return window extends past the as-of date.
  \item Fundamentals are joined only after their announcement dates, not by fiscal period alone.
  \item Feature selection and model tuning use only historical windows available before the simulated entry date in live-style validation.
  \item The Phase1 data snapshot ended at the April 30 A-share row; Phase2 used the newer May 8 row after those prices were public.
\end{itemize}

\clearpage

\section{Reproducibility And Acknowledgement}

Exact reproduction commands and hashes are documented in \texttt{README.md},
\texttt{docs/PHASE\_REPRODUCTION.md}, \texttt{docs/DATA\_PIPELINE.md},
\texttt{docs/MODEL\_SELECTION\_LINEAGE.md}, and
\texttt{docs/HYPERPARAMETERS.md}. The baseline comparison commands and outputs
are in \texttt{experiments/self\_test/SELF\_TEST\_REPORT.md}.

Instructor baseline files \texttt{baseline\_xgboost.py} and \texttt{features.py}
are retained unchanged for comparison. LLM assistance was used for coding,
documentation, and brainstorming, but the submitted stock weights were produced
by the student-trained XGBoost and portfolio-construction pipeline, not by
prompting an LLM for stock picks.

\end{document}
